% !PRINTED 524
\weberpage{524}

Nach \S.\,142, 4. k\"onnen wir, wenn $\omega$ ein gegebenes ganzes Functional, $E$ eine Einheit, $\varphi$ eine ganze Function in $\Omega$ ist, setzen:

\weq{1}{E\omega = \varphi(x, y, \ldots).}

Wenn wir andererseits $\omega$ in seine Primfactoren zerlegen, so k\"onnen wir nach 1. diese Primfactoren von den Variablen $x, y, \ldots$ frei annehmen, und wenn wir wieder das Product dieser Primfactoren bilden, so erhalten wir eine mit $\omega$ associirte Zahl $\omega_1$, die von den Variablen $x, y, \mathinner{.\,.}$ frei ist.

Ordnen wir den Quotienten $\varphi : \omega_1$, der eine ganze Zahl (sogar eine Einheit) ist, nach den Variablen $x, y, \ldots$, so schliessen wir aus \S.\,142, 3., dass alle Co\"efficienten der Function $\varphi$ durch $\omega_1$ und folglich auch durch $\omega$ und $\varphi$ theilbar sind.

Wenn andererseits alle Co\"efficienten von $\varphi$ durch irgend einen Factor $\delta$ in $\Omega$ theilbar sind, so ist auch $\varphi$ durch $\delta$ theilbar, und daraus ergiebt sich:

2. \gesperrt{Eine ganze Function $\varphi$ mit ganzen Zahlen als Co\"efficienten ist der gr\"osste gemeinschaftliche Theiler aller ihrer Co\"efficienten.}

Die ganze Zahl $\varphi$ geht also nach \S.\,139 in eine associirte Zahl \"uber, wenn die einzelnen Potenzen und Producte der Variablen durch je eine Variable ersetzt werden. Daraus ergiebt sich auch als Corollar, dass jedes Functional $\omega$ in ein associirtes \"ubergeht, wenn die Variablen irgendwie anders bezeichnet werden.

Es ist nun der folgende wichtige Satz zu beweisen:

3. \gesperrt{Ist $\omega$ ein beliebiges ganzes Functional, so kann man eine durch $\omega$ theilbare Zahl $\alpha$ so w\"ahlen, dass der Quotient $\alpha : \omega$ zu einem beliebig gegebenen Functionale $\mu$ relativ prim ist.}

Wir beweisen zun\"achst, dass es eine ganze Zahl $\alpha$ giebt, die durch $\omega$, aber nicht durch $\omega\pi$ theilbar ist, wenn $\pi$ ein beliebiges Primfunctional ist.

Bilden wir n\"amlich nach \S.\,142, 4. eine mit $\omega$ associirte ganze Function $\varphi$, so sind die Co\"efficienten dieser Function zwar alle durch $\omega$, aber nicht alle durch $\omega\pi$ theilbar, weil sonst auch $\varphi$ und mithin $\omega$ selbst durch $\omega\pi$ theilbar w\"are, was nicht m\"oglich ist. Es giebt also unter den Co\"efficienten von $\varphi$ wenigstens einen, der die verlangte Eigenschaft hat.
